% fixture: amsmath/align-label
% tier: T3
% engine: document
% covers: \label in align (\label@in@display, written by \make@display@tag in the tag's box after the number), in a first, a middle and a last row, beside \notag rows, with \ref and \eqref to them before and after
\documentclass{article}
\usepackage{amsmath}
\begin{document}
We refer ahead to \eqref{eq:b} and to \ref{eq:c}, which follow from
\begin{align}
a &= b + c \label{eq:a} \\
  &= d + e \notag \\
\label{eq:b} x^2 + y^2 &= z^2 \\
f(x) &= \int_0^x g(t)\,dt \nonumber \\
  &\le \frac{1}{2} \label{eq:c}
\end{align}
and back to \eqref{eq:a}, \eqref{eq:b} and \eqref{eq:c} in the text after the display.
\end{document}
